\chapter{Chief series of finite groups}
\label{ch:gr02-chief}

The composition series in D07 have simple factors and require
normality of each term in its successor
\cite[entries K07.06.composition\_series\_definition and
K07.06.jordan\_holder\_extension]{d07}. Here we refine the normal
subgroups of a fixed ambient group \(G\). Their successive quotients
carry the conjugation action of \(G\), which the factor comparison
will preserve.

\section{Chief factors and conjugation}

\begin{definition}[GR02-C02.01. Chief series and chief factors]
\label{def:gr02-chief-series}
Let \(G\) be a finite group. A chief series is a strictly increasing
chain
\[
  1=N_0<N_1<\cdots<N_r=G
\]
of normal subgroups of \(G\) such that no normal subgroup of \(G\)
lies strictly between \(N_{i-1}\) and \(N_i\), for \(1\le i\le r\).
The quotients \(N_i/N_{i-1}\) are the chief factors, and the integer
\(r\) is the length. The chain of length zero is allowed for \(G=1\).

A \(G\)-group is a group equipped with an action of \(G\) by
automorphisms. When \(K\le L\) are both normal in \(G\), the
natural \(G\)-group structure on \(L/K\) is
\[
  g\cdot(xK)=(gxg^{-1})K \qquad (g\in G,\ x\in L).
\]
A \(G\)-isomorphism of two such factors is a group isomorphism
which intertwines the two actions of the same ambient \(G\).
\end{definition}

The formula is well defined since conjugation by \(g\) preserves both
\(K\) and \(L\). The action restricted to \(L\) induces every inner
automorphism of \(L/K\), so a \(G\)-invariant subgroup is normal in
\(L/K\). For the trivial group, the list of factors is empty.

\begin{proposition}[GR02-C02.02. Existence and the minimal normal criterion]
\label{prop:gr02-chief-existence}
Every finite group has a chief series. A strictly increasing chain as
in Definition~\ref{def:gr02-chief-series}, with all its terms normal in
\(G\), is a chief series precisely when \(N_i/N_{i-1}\) is a minimal
nontrivial normal subgroup of \(G/N_{i-1}\) for each \(i\).
\end{proposition}
\begin{proof}
The finite poset of normal subgroups of \(G\) has a maximal chain
from \(1\) to \(G\). Such a chain admits no refinement by a normal
subgroup of \(G\). Normal subgroups of \(G/N_{i-1}\) contained in
\(N_i/N_{i-1}\) correspond under the quotient map to normal
subgroups of \(G\) between \(N_{i-1}\) and \(N_i\)
\cite[entry K07.04.correspondence\_theorem]{d07}. The endpoints of
the interval correspond to the trivial subgroup and the whole
factor. This proves the criterion in both directions.
\end{proof}

\section{Uniqueness and the solvable case}

\begin{theorem}[GR02-C02.03. Equivariant Jordan--H\"older uniqueness for chief factors]
\label{thm:gr02-chief-uniqueness}
Any two chief series of a finite group \(G\) have the same length
and the same multiset of \(G\)-isomorphism classes of factors.
\end{theorem}
\begin{proof}
Keep \(G\) fixed throughout. For every normal subgroup \(N\) of
\(G\), consider saturated chains from \(1\) to \(N\) in the poset
of normal subgroups of \(G\). Here saturated means that consecutive
terms admit no intermediate term in this poset. Such chains exist
by finiteness, and their factors all carry the \(G\)-action of
Definition~\ref{def:gr02-chief-series}.

We compare the multisets of factors of these chains by induction on
\(|N|\). The case \(N=1\) has an empty multiset. Let \(A\) and \(B\)
be the penultimate terms of two chains ending at a nontrivial
\(N\). If \(A=B\), induction compares the two chains up to \(A\)
and the top factor \(N/A\) is common to both. If \(A\ne B\), neither
contains the other by maximality. The product \(AB\) is normal in
\(G\) and strictly contains \(A\), so \(AB=N\). Put \(L=A\cap B\).
The quotient maps give isomorphisms
\[
  A/L \longrightarrow N/B,\quad aL\longmapsto aB,
  \qquad
  B/L \longrightarrow N/A,\quad bL\longmapsto bA.
\]
In each case the indicated intersection is the kernel before
passing to the quotient, and \(AB=N\) gives surjectivity.
Conjugating representatives by any \(g\in G\) commutes with both
maps, so these are \(G\)-isomorphisms.

There is no proper nontrivial \(G\)-invariant subgroup of \(N/B\),
since its inverse image in \(N\) would be normal in \(G\) and
strictly between \(B\) and \(N\). The same holds for \(N/A\).
By the two isomorphisms, \(L<A\) and \(L<B\) are therefore saturated
steps in the poset. Extend one saturated chain from \(1\) to \(L\)
first by \(A\) and then by \(N\), or first by \(B\) and then by \(N\).
Induction within \(A\) and \(B\) compares the original chains with
these two chains, including the common top step on each side.
The two constructed chains share all factors through \(L\), and
their last two factors match in reverse order by the displayed
\(G\)-isomorphisms. This finishes the induction. Taking \(N=G\)
proves the theorem.
\end{proof}

This is the descending comparison used for ordinary composition
series in D07 with the ambient conjugation retained at every step
\cite[entry K07.06.jordan\_holder\_extension]{d07}. The ordinary
group-isomorphism comparison for chief factors is also stated in
\cite[Theorem~5.9]{chief-source}.

\begin{corollary}[GR02-C02.04. Chief factors of finite solvable groups]
\label{cor:gr02-chief-solvable}
A finite group \(G\) is solvable if and only if it has a chief
series with elementary abelian factors. If \(G\) is solvable, every
chief factor is the additive group of a nonzero finite-dimensional
\(\mathbb F_p\)-vector space for some prime \(p\) depending on the
factor, and its \(G\)-action is linear over that field.
\end{corollary}
\begin{proof}
Assume \(G\) solvable. Passage to subgroups and quotients preserves
solvability by \cite[entry K07.15.solvable\_closure]{d07}. Hence the
chief factor \(D=N_i/N_{i-1}\) is solvable. Its derived subgroup
\([D,D]\) is characteristic in \(D\), hence is preserved by
conjugation by \(G\). Since \(D\) is nontrivial and solvable,
\([D,D]\) is proper. Minimality forces \([D,D]=1\), so \(D\) is
abelian. Choose a prime \(p\) dividing the order of some nonidentity
element of \(D\). A power of that element has order \(p\), and
\[
  D[p]=\{x\in D : x^p=1\}
\]
is a nontrivial characteristic subgroup because \(D\) is abelian.
Minimality now gives \(D[p]=D\). An abelian group of exponent \(p\)
has its usual \(\mathbb F_p\)-vector space structure by repeated
addition, and each group automorphism respects this scalar action.

Conversely, if the factors of a chief series are elementary
abelian then they are solvable. Closure under extensions makes
each term in the series solvable, starting at \(1\). This includes
the chain of length zero for \(G=1\).
\end{proof}

\begin{example}[GR02-C02.05. The two kinds of series in the alternating group]
\label{ex:gr02-chief-a4}
In \(A_4\), let
\[
  V=\{1,(12)(34),(13)(24),(14)(23)\},\qquad
  U=\langle(12)(34)\rangle .
\]
The subgroup \(V\) is elementary abelian of order \(4\), since the
product of two distinct double transpositions in \(V\) is the
third. It is normal in \(A_4\) since conjugation preserves cycle
type. Conjugation by \((123)\) permutes the three nonidentity
elements of \(V\) cyclically. Consequently every nontrivial
normal subgroup of \(A_4\) contained in \(V\) equals \(V\).
As \(A_4/V\) has order \(3\), \(1<V<A_4\) is a chief series.
Its factors are \(V\) and \(C_3\). The action on \(V\) factors
through \(A_4/V\) and permutes its three nonzero vectors, while
the conjugation action on the factor \(A_4/V\) is trivial.

The longer chain \(1<U<V<A_4\) is a composition series since
adjacent terms are normal and the factors have prime orders
\(2,2,3\). The subgroup \(U\) is moved to another subgroup by
\((123)\), so this chain fails the ambient normality condition
in Definition~\ref{def:gr02-chief-series}. See also
\cite[Examples~5.2 and 5.7]{chief-source}.
\end{example}
